\documentclass{article}

% ------------------------------------------------------------ page style
% The arXiv preprint look of arxiv.sty (github.com/kourgeorge/arxiv-style).
% That package is not part of TeX Live, so Overleaf cannot find it. The parts
% of it this document uses are set here, and the file compiles on its own.
\makeatletter
\renewcommand{\rmdefault}{ptm}
\renewcommand{\sfdefault}{phv}
\usepackage[letterpaper]{geometry}
\AtBeginDocument{\newgeometry{textheight=9in, textwidth=6.5in, top=1in,
  headheight=14pt, headsep=25pt, footskip=30pt}}
\widowpenalty=10000
\clubpenalty=10000
% Ragged rather than flush bottom: tables never break, so a page that has to end
% early keeps its white space at the foot instead of spreading it between lines.
\raggedbottom
\sloppy

\newcommand{\shorttitle}{\@title}
\usepackage{fancyhdr}
\fancyhf{}
\pagestyle{fancy}
\renewcommand{\headrulewidth}{0.4pt}
\fancyheadoffset{0pt}
\chead{\shorttitle}
\cfoot{\thepage}

% font sizes with reduced leading
\renewcommand{\normalsize}{%
  \@setfontsize\normalsize\@xpt\@xipt
  \abovedisplayskip      7\p@ \@plus 2\p@ \@minus 5\p@
  \abovedisplayshortskip \z@ \@plus 3\p@
  \belowdisplayskip      \abovedisplayskip
  \belowdisplayshortskip 4\p@ \@plus 3\p@ \@minus 3\p@}
\normalsize
\renewcommand{\small}{%
  \@setfontsize\small\@ixpt\@xpt
  \abovedisplayskip      6\p@ \@plus 1.5\p@ \@minus 4\p@
  \abovedisplayshortskip \z@  \@plus 2\p@
  \belowdisplayskip      \abovedisplayskip
  \belowdisplayshortskip 3\p@ \@plus 2\p@   \@minus 2\p@}
\renewcommand{\footnotesize}{\@setfontsize\footnotesize\@ixpt\@xpt}
\renewcommand{\large}{\@setfontsize\large\@xiipt{14}}
\renewcommand{\LARGE}{\@setfontsize\LARGE\@xviipt{20}}
\renewcommand{\huge}{\@setfontsize\huge\@xxpt{23}}

% headings with less space
\renewcommand{\section}{\@startsection{section}{1}{\z@}%
  {-2.0ex \@plus -0.5ex \@minus -0.2ex}{1.5ex \@plus 0.3ex \@minus 0.2ex}%
  {\normalfont\large\bfseries\raggedright}}
\renewcommand{\subsection}{\@startsection{subsection}{2}{\z@}%
  {-1.8ex \@plus -0.5ex \@minus -0.2ex}{0.8ex \@plus 0.2ex}%
  {\normalfont\normalsize\bfseries\raggedright}}
\renewcommand{\paragraph}{\@startsection{paragraph}{4}{\z@}%
  {1.5ex \@plus 0.5ex \@minus 0.2ex}{-1em}%
  {\normalfont\normalsize\bfseries}}

% paragraphs, and the list spacing that center and quote inherit
\setlength{\parindent}{\z@}
\setlength{\parskip}{5.5\p@}
\setlength{\topsep}{4\p@ \@plus 1\p@ \@minus 2\p@}
\setlength{\partopsep}{1\p@ \@plus 0.5\p@ \@minus 0.5\p@}
\setlength{\itemsep}{2\p@ \@plus 1\p@ \@minus 0.5\p@}
\setlength{\parsep}{2\p@ \@plus 1\p@ \@minus 0.5\p@}
\setlength{\leftmargin}{3pc}
\setlength{\leftmargini}{\leftmargin}
\def\@listi{\leftmargin\leftmargini}

% title block and abstract
\renewcommand{\maketitle}{\par
  \begingroup
    \thispagestyle{empty}%
    \vbox{\hsize\textwidth \linewidth\hsize
      \vskip 0.1in
      \hrule height 2\p@ \vskip 0.25in \vskip -\parskip
      \centering
      {\LARGE\scshape \@title\par}%
      \vskip 0.29in \vskip -\parskip \hrule height 2\p@
      \vskip 0.3in}%
  \endgroup
  \let\maketitle\relax}
\renewenvironment{abstract}
  {\centerline{\large\bfseries\scshape Abstract}\begin{quote}}
  {\end{quote}}
\makeatother

\usepackage[utf8]{inputenc}
\usepackage[T1]{fontenc}
\usepackage{hyperref}
\usepackage{url}
\usepackage{booktabs}
\usepackage{amsmath}
\usepackage{amssymb}
\usepackage{graphicx}
\usepackage{xcolor}
\usepackage{listings}
\usepackage{array}
\usepackage{tabularx}
\usepackage{needspace}

\hypersetup{colorlinks=true, linkcolor=black, citecolor=black, urlcolor=blue}

% ---------------------------------------------------------------- macros
% Every section carries the script that produced its numbers and the span
% they came from. \prov is that line. It is treated as part of the heading, so
% a page cannot end right after it and strand the heading at the foot.
\makeatletter
\newcommand{\prov}[2]{%
  \par\vspace{1pt}\noindent
  {\small\texttt{#1}\ \ (\textit{#2})}\par
  \nobreak\vspace{5pt}\@afterheading}
\makeatother

% The forward pointer: what this section's conclusion obliges the next one to
% measure. The chain of these is the argument of the paper.
% amssymb already defines \leadsto (a math arrow, unused here), so renew it.
\renewcommand{\leadsto}[1]{%
  \par\vspace{3pt}\noindent
  \textcolor{black!70}{$\blacktriangleright$}\ \textbf{Leads to.}\ #1\par}

\newcommand{\bic}{\texttt{block\_ic}}
\newcommand{\vsbh}{\textit{vs bh}}
\newcommand{\detec}{\texttt{detectable}}

\lstset{basicstyle=\ttfamily\footnotesize, frame=single, framesep=4pt,
        xleftmargin=6pt, breaklines=true, columns=flexible}

\title{Predict, then size: deriving a daily position rule from a signal that
       orders returns but cannot scale them}

\renewcommand{\shorttitle}{Predict, then size}

\begin{document}
\maketitle

This report records how a daily position rule was derived, one step at a time. It
covers three things: how the features were built, how they and the model's output
were diagnosed, and how the resulting prediction becomes a position.

Each step names the script that produced its numbers and the span they came from,
gives the result, and says what the result forces the next step to measure. Reading
only those hand-offs gives the whole argument. Code and cached output sit beside
this file in the repository.

\section*{Setup}
\label{sec:setup}

\paragraph{The data.}
Daily S\&P 500 forward returns, the risk-free rate, and predictors named
\texttt{<group><number>}, from the Kaggle \textit{Hull Tactical --- Market
Prediction} series. The competition documents eight groups; this extract carries
93 raw columns across seven of them.

\begin{center}\small
\begin{tabular}{llr}
\toprule
group & meaning & raw columns \\
\midrule
\texttt{E} & economic and macro indicators & 19 \\
\texttt{M} & market and price-based indicators & 18 \\
\texttt{P} & positioning and valuation & 13 \\
\texttt{V} & volatility measures & 13 \\
\texttt{S} & sentiment signals & 12 \\
\texttt{D} & date and calendar flags & 9 \\
\texttt{I} & interest-rate and institutional series & 9 \\
\texttt{MOM} & momentum factors & --- \\
\bottomrule
\end{tabular}
\end{center}

\noindent The prefix \texttt{MOM} is matched before \texttt{M}, so a momentum
series would never fall into the market group. Columns start at different dates,
and the raw series runs 8990 rows against the feature table's 8918.

\paragraph{The split.}
\texttt{dev} (69 blocks, about 23 years, rows 1261--7056) fits models and measures
mechanisms. \texttt{valid} (10 blocks, rows 7141--7980) chooses between
configurations. \texttt{test} (11 blocks, rows 7981--8904) is read once, at the
end. The development boundary is row 7113; the first holdout block starts at 7141.
Part~1 touches neither holdout span.

\paragraph{The metric.}
The competition's adjusted Sharpe. A position $w \in [0,2]$ earns
$r_f(1-w) + w\,f_r$, and the Sharpe of that series is divided by
\[
  \Big(1 + \max\big(0,\ \sigma_s/\sigma_m - 1.2\big)\Big)
  \times
  \Big(1 + \mathrm{gap}^2/100\Big),
\]
where $\sigma_s/\sigma_m$ is realised strategy volatility over market volatility
and $\mathrm{gap}$ is the shortfall of mean return against the market.
Buy-and-hold is $w=1$ and is also the market, so its adjusted Sharpe equals its
plain Sharpe. For every configuration reported here both penalty factors come out
at exactly 1.000, so the adjusted Sharpe below is the plain Sharpe.

\paragraph{Terms.}
\bic{} is the mean per-block Spearman correlation between the signal and the next
day's excess return, one block being the 84 rows a single fit predicts.
\texttt{spread\_ann} is the annualised top-minus-bottom-decile return, averaged
over folds. \detec{} is twice the standard error of a paired contrast; a
difference smaller than it cannot be separated from zero on the folds available.
It is reported beside every contrast in place of a $p$-value.

\paragraph{Naming.}
The scripts that sweep the same parameter name it differently. This report uses
\texttt{pipeline.py}'s names.

\begin{center}\small
\begin{tabular}{lllll}
\toprule
meaning & \texttt{pipeline} & \texttt{mapping\_form} & \texttt{detrend\_sweep} A/B & \texttt{detrend\_sweep} C \\
\midrule
rolling mean removed before ranking & \texttt{detrend} & \texttt{mu\_window} & \texttt{roll21} & \texttt{mean8}, \texttt{ewm5} \\
the percentile's reference window & \texttt{rank\_window} & \texttt{rank\_window} & \texttt{rank\_win} & \texttt{rank\_win} \\
EWMA half-life on the position & \texttt{smooth\_halflife} & --- & --- & \texttt{smooth5} \\
\bottomrule
\end{tabular}
\end{center}

So \texttt{mean8} $\times$ \texttt{smooth5} $\times$ \texttt{rank252} reads:
subtract the mean of the last 8 predictions, take the percentile within the last
252 days, solve the scale, then apply an EWMA of half-life 5 days to the position.

\section{The signal layer}

\subsection{Data processing}
\label{sec:features}
\prov{Hull\_Tactical\_feature\_engineering.py}{the whole table}

93 raw columns become 1132 features. What follows is what is built, on what
window, and whether any of it can see the future.

\textit{Coverage.} A column is kept when its data begins in the first 70\% of the
series, which admits late-starting predictors rather than truncating the panel to
the shortest one. The per-group quantile rule is present and disabled.

\textit{Lagged returns.} The three return series are shifted by one row before
anything is built on them, so a signal at $t$ uses returns through $t-1$. They
carry their own rolling block at ultra-short horizons (2, 3, 5, 10, 21), 74
columns.

\textit{Rolling statistics.} Each raw column gets a rolling mean, z-score and
range position at every window its group is assigned, plus one first difference.
The windows are group-specific, because a volatility series and a macro series
move on different clocks.

\begin{center}\small
\begin{tabular}{ll}
\toprule
group & windows \\
\midrule
\texttt{V} & 5, 21, 63 \\
\texttt{S} & 5, 10, 21 \\
\texttt{M} & 5, 10, 21, 63 \\
\texttt{MOM} & 5, 21 \\
\texttt{E}, \texttt{P} & 63, 126, 252 \\
\texttt{I} & 21, 63 \\
\texttt{D} & 63, rolling mean only \\
\bottomrule
\end{tabular}
\end{center}

\noindent \texttt{D} columns are 0/1 flags, so only the raw value and a rolling
mean survive, which records how often the flag was on recently. The \texttt{\_mean}
columns of \texttt{P}, \texttt{I} and \texttt{S} are dropped once the gates are
built, since the level of a slow series adds little beside its z-score and range
position. The result is 289 z-scores, 267 range positions, 199 rolling means and
88 first differences.

\textit{Group PCA, refitted every row.} Each group's z-scored columns are reduced
to a first principal component. At row $t$ the component is fitted on the window
$[t - w,\, t)$ and applied to row $t$ alone, so the fit never sees the row it
transforms. Only columns with data somewhere in that window take part, which lets
the active set grow as late-starting series come online, and the inputs are
winsorised at the window's own 98th percentile before the fit.

\begin{center}\small
\begin{tabular}{lllll}
\toprule
group & method & window & warm-up & z-window / input \\
\midrule
\texttt{V} & rolling & 126 & 126 & 126 / \texttt{\_z21} \\
\texttt{M} & rolling & 126 & 126 & 126 / \texttt{\_z21} \\
\texttt{MOM} & rolling & 63 & 63 & 126 / \texttt{\_z21} \\
\texttt{E} & rolling & 252 & 252 & 252 / \texttt{\_z63} \\
\texttt{S} & rolling & 126 & 126 & 252 / \texttt{\_z21} \\
\texttt{I} & expanding & --- & 504 & 252 / \texttt{\_z21} \\
\texttt{P} & expanding & --- & 504 & 252 / \texttt{\_z63} \\
\bottomrule
\end{tabular}
\end{center}

\noindent \texttt{I} and \texttt{P} expand rather than roll because their series
are slow. Each group produces \texttt{pc1}, its z-score and its positive part, 18
columns over the six groups with enough active inputs.

\textit{Regime gates.} The PC1 factors are the regime variables:
\[
  \text{gate} \;=\; 1 + \mathrm{clip}(\texttt{pc1\_z},\, 0,\, 3) \;\in\; [1, 4],
\]
neutral at 1. Selected z-score columns are multiplied by it, so a feature counts
for more when its own group's regime is extreme and is left alone otherwise. The
gated windows are short for the fast groups and long for the slow ones:
\texttt{V}, \texttt{S}, \texttt{M} and \texttt{MOM} at 21 and 63, \texttt{E} and
\texttt{P} at 126 and 252, \texttt{I} at 63 and 126. That adds 159 columns. Five
hand-built interactions are also kept under an \texttt{XINT\_} prefix:
$\texttt{I2}-\texttt{I1}$, $\texttt{I7}-\texttt{I1}$, $\texttt{P1}/\texttt{E1}$,
$\texttt{V1}/\texttt{M1}$, $\texttt{S1}-\texttt{S2}$.

\textit{Two assumptions worth naming.} \texttt{E} and \texttt{S} columns are
shifted by 10 rows before any window is taken, on the assumption that news in
those blocks reaches prices with a delay. That is an assumption, not a
measurement. And nothing is pruned: \S\ref{sec:ablation} removes each structural
block in turn and every removal costs or ties, so all 1132 columns stay.

\medskip
\noindent Now the causality question. Every dispersion statistic is computed on
\texttt{s.rolling(w, min\_periods=w)}, and the group factors come from
\texttt{causal\_pc1\_onfly\_winsor} above. The builder contains no
\texttt{fit\_transform} and no global mean or standard deviation, so a feature at
row $t$ uses only rows $\le t$.

One dependency on the whole series survives. The coverage filter measures how
late a column starts as a fraction of the table's full 8918 rows, and drops it
below 0.30. Normalising by the whole table rather than by development alone puts
the threshold at row 6243 instead of 4979, and the stricter version would drop 71
more columns. The non-causal input is a scalar, ``how long is the dataset''. No
holdout value participates, and the latest first-valid row among surviving columns
is 6039, well before the boundary at 7113.

\subsection*{The split, cut once}
\label{sec:split}
\prov{split\_data.py}{after \S\ref{sec:features}, before everything else}

The feature table is cut into three files, and every later script reads a span
through \texttt{paths.load\_split(span)}. The geometry is derived from
\texttt{wfo.make\_folds} under the frozen constants rather than written down, so it
cannot drift from what the models do.

\begin{center}\small
\begin{tabular}{lrrrr}
\toprule
 & rows kept & scored & folds & history \\
\midrule
\texttt{train} & 504--7112 & 1261--7056 & 69 & 757 \\
\texttt{valid} & 6384--7980 & 7141--7980 & 10 & 757 \\
\texttt{test}  & 7224--8917 & 7981--8904 & 11 & 757 \\
\bottomrule
\end{tabular}
\end{center}

A span's models train on rows that precede it, so each file carries its scored
rows plus the 757 rows of history those fits need: the training window, the purge
and the embargo. The files therefore overlap, and any one of them runs the full
walk-forward for its span alone. Each file is an exact slice, so
\texttt{feature\_cols} still returns 1132 and \texttt{manifest.json} records where
each slice begins.

Three entry points cover the later uses. \texttt{load\_span} returns a slice and
its folds, indexed into it. \texttt{load\_through} returns every row from the first
training row through the end of a span, which one fit on all of development needs
because it trains on 6636 rows while valid's own slice carries 757.
\texttt{remap\_rows} and \texttt{slice\_covering} place a script's own fold grid
onto a slice and refuse, loudly, when the slice lacks the history that grid needs.

The split runs after the builder because the builder drops the leading warm-up
rows, \texttt{date\_id} 0 to 71, where the longest rolling windows have no history.
The raw series has 8990 rows and the feature table 8918, and the fold grid is
positional on the feature table, so a split computed on the raw index lands 58 rows
away from where the models cut.

\texttt{test} is now unreachable unless a script names it, and nothing before the
final scoring does. That closes two failure modes: a holdout row reaching a fit
through a trailing statistic, and a fold grid quietly losing folds because a slice
was too short to hold them.

\subsection{Check that the group factors are still factors}
\label{sec:pc1}
\prov{pc1\_check.py}{train split}

\texttt{MIN\_COVERAGE} $= 0.30$ lets late-starting columns back in, so each group's
PCA active set now grows several times. Counting sign flips and jumps in PC1 as
columns come online puts both at about 4--5\% of rows. A rolling z-score absorbs a
scale change but not a sign change, so that rate is the thing to watch.

Small enough to keep the factors, large enough to record:
\texttt{V\_pc1\_z126} and \texttt{E\_pc1\_z252} are downstream of it.

\leadsto{The features are usable. Before asking which of them matter, the model
that answers has to be large enough to reach them.}

\subsection{Fix the model's capacity first}
\label{sec:capacity}
\prov{capacity\_window\_grid.py}{dev, 69 folds, causal trailing percentile}

Training window against leaf count, as a $3 \times 5$ grid rather than two ladders,
because the axes interact: the capacity gain from 5 to 62 leaves measured $+0.442$
at a 756-row window and $+0.124$ at an expanding one. Scored by annualised decile
spread, with paired contrasts on shared folds.

\begin{center}\small
\begin{tabular}{lrrrrr}
\toprule
window $\backslash$ target & 5 & 20 & 62 & 117 & 176 \\
\midrule
252 & 0.183 & 0.519 & 0.670 & --- & --- \\
504 & 0.446 & 0.672 & 0.672 & 0.626 & --- \\
\textbf{756} & 0.364 & 0.699 & 0.704 & \textbf{0.854} & 0.928 \\
\bottomrule
\end{tabular}
\end{center}

\begin{center}\small
\begin{tabular}{lrrrl}
\toprule
contrast & diff & $t$ & \detec{} & separable \\
\midrule
@756: T62 $-$ T5    & $+0.511$ & 2.11 & 0.507 & yes \\
@756: T117 $-$ T5   & $+0.669$ & 2.52 & 0.555 & yes \\
@756: T176 $-$ T5   & $+0.737$ & 2.84 & 0.545 & yes \\
@T5: W252 $-$ W756  & $-0.171$ & $-0.76$ & 0.464 & no \\
@T20: W252 $-$ W756 & $-0.247$ & $-1.11$ & 0.458 & no \\
@T62: W252 $-$ W756 & $-0.180$ & $-0.79$ & 0.470 & no \\
@T117: W504 $-$ W756& $-0.039$ & $-0.17$ & 0.491 & no \\
W756T176 $-$ W756T62& $+0.195$ & 1.85 & 0.215 & no \\
\bottomrule
\end{tabular}
\end{center}

Capacity separates and the window does not. Above 62 leaves the gains stop
separating from each other, so 62, 117 and 176 form one band rather than a ranking.
Frozen at window 756, target 117: the middle of that band, at the longest window.
\texttt{embargo\_capacity.py} later re-swept $\{5,20,62,117\}$ at embargo 0 and put
117 top on dev again, \bic{} 0.1011 against 0.0943 at 62.

\leadsto{The model is now large enough that an ablation reports on the features
rather than on the model. \S\ref{sec:ablation} shows what the reversed order
produces.}

\subsection{Ablate the feature blocks}
\label{sec:ablation}
\prov{ablate\_blocks.py}{dev, 69 folds, window 756, target 117, causal trailing percentile}

Remove one structural block at a time, at a capacity the model can use. On a
756-row window:

\begin{center}\small
\begin{tabular}{lrrr}
\toprule
 & splits/tree & distinct columns/tree & columns ever used \\
\midrule
4 leaves  &  3 &  2.9 & \textbf{267 of 1132 (24\%)} \\
92 leaves & 91 & 65.7 & 666 of 1132 (59\%) \\
\bottomrule
\end{tabular}
\end{center}

A model that touches a quarter of the columns cannot report on the other three
quarters. This runs at 117.

\begin{center}\small
\begin{tabular}{lrrrrrr}
\toprule
removed & columns & \texttt{spread\_ann} & vs baseline & $t$ & \detec{} & at 4 leaves \\
\midrule
baseline & 0 & 0.854 & --- & --- & --- & --- \\
\textbf{D group} & 18 & 0.707 & $\mathbf{-0.187}$ & $\mathbf{-2.81}$ & 0.136 & $-0.067$ \\
long \texttt{\_mean} & 97 & 0.754 & $-0.095$ & $-1.32$ & 0.147 & $+0.013$ \\
all \texttt{\_pos} & 261 & 0.691 & $-0.172$ & $-1.78$ & 0.197 & $+0.024$ \\
long \texttt{\_z} & 64 & 0.864 & $+0.073$ & $+0.98$ & 0.152 & $+0.037$ \\
gates & 159 & 0.699 & $-0.109$ & $-1.32$ & 0.169 & $-0.051$ \\
\bottomrule
\end{tabular}
\end{center}

Removing any block costs or ties, and nothing gains. Only the D group separates,
and it is a loss, so all 1132 columns stay. The last column is the same
measurement at four leaves, where every effect sits inside $\pm 0.067$: a uniform
null that says nothing about the features and everything about the model producing
it.

\leadsto{Capacity and features are fixed. What remains of the model is the
hyperparameters underneath them.}

\subsection{The remaining tree hyperparameters}
\label{sec:hyper}
\prov{tree\_sensitivity.py}{dev, 69 folds, one axis at a time, within-block ranking}

\begin{center}\small
\begin{tabular}{lrrrr}
\toprule
axis & default & best alternative & difference & \detec{} \\
\midrule
\texttt{learning\_rate}    & \textbf{0.05} & 0.01 & $-0.060$ & 0.198 \\
\texttt{colsample\_bytree} & \textbf{0.7}  & 0.9  & $-0.065$ & 0.242 \\
\texttt{reg\_lambda}       & \textbf{5.0}  & 1.0  & $-0.051$ & 0.210 \\
\texttt{subsample}         & \textbf{0.7}  & 0.85 & $-0.176$ & 0.244 \\
\texttt{n\_estimators}     & \textbf{250}  & 500  & $+0.010$ & 0.065 \\
\bottomrule
\end{tabular}
\end{center}

Every deviation is negative or inside its own resolution, so the defaults stand.
The extremes are genuine losses: \texttt{learning\_rate} 0.20 at $-0.528$
($t = -4.22$), \texttt{reg\_lambda} 100 at $-0.651$, \texttt{colsample} 1.0 at
$-0.268$ ($t = -2.70$).

\leadsto{The model is fully specified. What decides the rest of the project is
which part of its output is usable: the number, or only the ordering.}

\subsection{What kind of signal this is}
\label{sec:signalkind}
\prov{ridge\_signal.py (both models), signal\_diagnostics.py (LightGBM)}{dev, 69 folds}

Three measurements. The level: $R^2_{\text{oos}}$ against zero and against the
trailing mean, then the Mincer--Zarnowitz slope with Newey--West errors, then that
slope re-estimated from earlier rows only and applied forward, which tests whether
the overstatement was knowable at the time. The order: per-block Spearman. The
increment: inside each 84-row block, $r$ is the normal score of the within-block
rank and $m = (s - \text{block mean})/\text{block sd}$, so $r$ carries the order
alone and $m$ carries order and magnitude together.

\begin{center}\small
\begin{tabular}{lrrrr}
\toprule
level & sd(pred) & sd(truth) & $R^2$ vs zero & vs trailing mean \\
\midrule
\textbf{ridge} & 0.0062 & 0.0110 & $\mathbf{-0.261}$ & $-0.260$ \\
tree           & 0.0044 & 0.0110 & $-0.119$ & $-0.117$ \\
\bottomrule
\end{tabular}
\end{center}

\begin{center}\small
\begin{tabular}{lrrrrr}
\toprule
 & MZ slope $b$ & NW $t$ & $1/b$ & causal $b$ & $R^2$ after causal rescale \\
\midrule
\textbf{ridge} & \textbf{0.0872} & \textbf{3.03} & \textbf{11.5$\times$} & 0.0899 & $\mathbf{+0.0015}$ \\
tree           & 0.1393 & 2.54 & 7.2$\times$ & 0.1431 & $+0.0015$ \\
\bottomrule
\end{tabular}
\end{center}

\begin{center}\small
\begin{tabular}{lrrrr}
\toprule
order & \bic{} & $t$ & blocks positive & pooled IC \\
\midrule
\textbf{ridge} & \textbf{0.1060} & \textbf{10.93} & \textbf{91.3\%} & 0.0448 \\
tree           & 0.1011 & 8.11 & 87.0\% & 0.0531 \\
\bottomrule
\end{tabular}
\end{center}

\begin{center}\small
\begin{tabular}{llrrr}
\toprule
increment & model & $R^2$ & $t$ on $r$ & $t$ on $m$ \\
\midrule
\textbf{ridge} & $r$ only & 0.0108 & 8.36 & \\
               & $m$ only & 0.0109 & & 8.28 \\
               & \textbf{both} & \textbf{0.0110} & \textbf{0.46} & \textbf{0.74} \\
tree & $r$ only & 0.0094 & 7.57 & \\
     & $m$ only & 0.0117 & & 8.10 \\
     & \textbf{both} & \textbf{0.0136} & $-2.41$ & $\mathbf{+3.52}$ \\
\bottomrule
\end{tabular}
\end{center}

A linear signal exists at $t = 3.03$, and the prediction overstates it by
$11.5\times$. The overstatement was knowable at the time, which is why $R^2$ moves
from $-0.261$ to $+0.0015$ once the slope is estimated causally and applied
forward. The order is strong. For the Ridge, adding the magnitude to the order buys
$+0.0002$ of $R^2$ and neither coefficient survives, so the magnitude can be
thrown away for free. Keep the rank and map it through a percentile. Every rule
that converts a prediction into return units before sizing is out.

For the tree, $m$ stays significant and joint $R^2$ rises from 0.0094 to 0.0136, so
a rank-based mapping handicaps it specifically. Part~2's search therefore owes the
tree a magnitude encoding; \S\ref{sec:ridge} settles that debt before dropping it.

\leadsto{The signal layer is closed and produces one number per day, an ordering.
Part~2 asks what position that ordering becomes, starting with the shape of the
curve.}

\section{The mapping layer}

\subsection{Every shape of the rank-to-position curve works}
\label{sec:curve}
\prov{mapping\_form.py -{}-curves}{dev, 69 folds, $\bar w = 1.0$, scale solved causally to a volatility ratio of 1.2}

Part~1 leaves one number per day, where today sits among the last 252. Seven curves
from the literature turn that into a position: AFML sigmoid sizing,
Gu--Kelly--Xiu extreme-decile spread, dead bands and the rest. Each gets its own
causally solved scale, so none can win by taking more risk. The shapes differ
substantially:

\begin{center}\small
\begin{tabular}{lrrrrr}
\toprule
$p$ & linear & normal score & decile & extreme decile & dead band \\
\midrule
0.02 & 0.31 & 0.21 & 0.34 & 0.31 & \textbf{0.02} \\
0.25 & 0.64 & 0.74 & 0.63 & \textbf{1.00} & \textbf{1.00} \\
0.50 & 1.00 & 1.00 & 1.07 & 1.00 & 1.00 \\
0.75 & 1.36 & 1.26 & 1.37 & \textbf{1.00} & \textbf{1.00} \\
0.98 & 1.69 & 1.79 & 1.66 & 1.69 & \textbf{1.98} \\
\bottomrule
\end{tabular}
\end{center}

\begin{center}\small
\begin{tabular}{lrrrrrrr}
\toprule
curve & pooled adj & vs market & $t$ & vs linear & \detec{} & turnover & days with a view \\
\midrule
normal score   & 0.6103 & $+0.2817$ & 3.22 & $+0.0019$ & 0.045 & 0.257 & 98.6\% \\
\textbf{linear}& 0.6015 & $+0.2799$ & 3.25 & --- & --- & 0.274 & 98.6\% \\
decile 10      & 0.5866 & $+0.2813$ & 3.36 & $\mathbf{+0.0014}$ & \textbf{0.021} & 0.271 & 100\% \\
multi-horizon  & 0.5862 & $+0.2942$ & 3.57 & $+0.0143$ & 0.050 & 0.288 & 99.0\% \\
extreme decile & 0.5807 & $+0.2820$ & 2.79 & $+0.0021$ & 0.143 & 0.190 & 21.5\% \\
dead band      & 0.5777 & $+0.2798$ & 3.24 & $-0.0000$ & 0.132 & \textbf{0.159} & \textbf{21.5\%} \\
isotonic       & 0.4788 & $+0.2032$ & 4.60 & $-0.0766$ & 0.128 & 0.118 & 95.7\% \\
\bottomrule
\end{tabular}
\end{center}

All seven beat buy-and-hold, by $+0.28$ per fold at $t$ between 2.8 and 4.6, and
pooled 0.60 against the market's 0.406. None separates from any other: the tightest
null is decile 10 at $+0.0014$ against a \detec{} of 0.021. Turnover varies by
$1.8\times$ and time in the market by $4.6\times$ with no change in score. Once the
risk budget is fixed, any monotone transform of the rank extracts the same
information and the shape only moves where the risk is taken. Take the linear curve
and stop tuning shape.

This also sets the baseline the rest of Part~2 has to explain. Carried forward
unchanged, this mapping scores $-0.390$ on valid, so the shape is not what fails
out of sample.

One residual. The normaliser divides by a full-sample standard deviation, but the
causal solver cancels it exactly: scaling $z$ by $c$ makes the solver return
$\tau/c$, verified at $4.4\times10^{-16}$ beyond row 120. The residual covers 120
of 5796 rows and is identical across all seven candidates. The isotonic fit is
refitted every 84 rows on rows strictly before the block, and no fold reaches past
row 7066.

\leadsto{The curve is fixed and its amplitude is still free. The metric has an
opinion about amplitude, so that comes next.}

\subsection{How much of the signal can be deployed}
\label{sec:amplitude}
\prov{mapping\_form.py -{}-params}{dev, 69 folds; LightGBM evidence, corroborated on the Ridge in \S\ref{sec:budget}}

The position is $w = 1 + \tau z$, with $\tau$ solved so $w$ runs at a chosen ratio
of market volatility. The metric taxes ratios above 1.2 linearly, so the question is
whether to stop at the kink. Reading the unpenalised Sharpe answers it, because the
tax is divided out and what is left measures the signal.

\begin{center}\small
\begin{tabular}{lrrrr}
\toprule
budget & adjusted & unpenalised Sharpe & vol penalty & mean $\tau$ \\
\midrule
1.0 & 0.410 & 0.410 & 1.000 & 0.009 \\
1.1 & 0.550 & 0.550 & 1.000 & 0.291 \\
\textbf{1.2} & \textbf{0.602} & 0.602 & 1.000 & 0.464 \\
1.3 & 0.572 & \textbf{0.617} & 1.080 & 0.609 \\
1.4 & 0.541 & \textbf{0.617} & 1.140 & 0.740 \\
1.5 & 0.524 & \textbf{0.616} & 1.176 & 0.862 \\
1.7 & 0.509 & \textbf{0.619} & 1.216 & 1.090 \\
2.0 & 0.495 & \textbf{0.617} & 1.246 & 1.409 \\
\bottomrule
\end{tabular}
\end{center}

\begin{center}\small
\begin{tabular}{lrrr}
\toprule
how $\tau$ is obtained & pooled adjusted & realised ratio & vol penalty \\
\midrule
one solve per fold, held 84 days & 0.6281 & \textbf{1.234} & \textbf{1.034} \\
\textbf{re-solved every day} & \textbf{0.6773} & 1.199 & 1.000 \\
\bottomrule
\end{tabular}
\end{center}

The unpenalised Sharpe climbs $+0.19$ between ratios 1.0 and 1.2, then sits at
0.617 across a doubling of the swing while the tax keeps growing. Past about 1.3 a
wider swing buys variance and no return. So $\bar w = 1.0$, budget 1.2, $\tau$
re-solved daily on a trailing 756 rows. Crossing the kink is out, and so is picking
$\tau$ by hand.

The daily re-solve is worth $+0.049$ pooled, the same order as the budget itself,
while its per-fold paired difference is $-0.0096$ against a \detec{} of 0.115. Only
the pooled number can see it, because the metric's penalties are computed on
whatever window they are handed. The solver's own constants move the result by less
than $\pm 0.005$: \texttt{tau\_min\_hist} $60/120/252 \to 0.6054/0.6015/0.6011$,
\texttt{tau\_fallback} $0.0/0.3/0.6/1.0 \to 0.6008/0.6015/0.6020/0.6037$, solver
window $252/504/756 \to 0.580/0.600/0.602$. As a check, budget 1.0 scores 0.410
with mean $\tau = 0.009$, reproducing the market's 0.406.

The saturation is a property of the signal and carries to any span. The value 1.2
is this metric's kink and carries nowhere.

\leadsto{Curve and amplitude are fixed. One degree of freedom is left, what the
percentile is measured against, and the next four steps are spent on it.}

\subsection{How much is there to gain from the reference set}
\label{sec:reference}
\prov{detrend\_sweep.py -{}-grid}{Ridge, dev, 69 folds}

The shipped rule ranks today against the trailing 252 days after removing a 63-day
rolling mean. The alternative reference is the 84 rows this fit produced, since each
model predicts exactly one block and block and fit are the same thing. Re-ranking
the same predictions against each reference prices the difference, and splitting the
within-block rank into the block's mean alone and then the full rank gives two
look-ahead ceilings.

\begin{center}\small
\begin{tabular}{lr}
\toprule
reference & \bic{} \\
\midrule
trailing 252, after a 63-day mean (shipped) & 0.0756 \\
trailing 252, no mean removed & \textbf{0.0950} \\
the block's actual mean removed (look-ahead) & 0.1002 \\
the full within-block rank (look-ahead) & 0.1060 \\
\bottomrule
\end{tabular}
\end{center}

Ranking against a trailing window with nothing removed lands within 0.0053 of
knowing the block's true centre, and within 0.0110 of the full within-block rank.
The Ridge therefore has almost no level problem. The shipped rule gives away 0.0247
against perfect centring, and 0.0194 of that, four fifths, is the 63-day mean
damaging the signal rather than a level the reference set fails to remove. For this
model, 63 is the wrong length.

That also closes the whole line of work that estimates a fit's output level and
subtracts it. A perfect level estimate is worth 0.0053 here, and
\texttt{level\_refit.py -{}-model ridge} prices the best available one directly: it
correlates at $+0.342$, which is 2.9 standard errors and therefore genuinely not
flat, but its dispersion is a fifth of the truth's, 0.0009 against 0.0042. It
captures roughly 12\% of the level variance and buys about $+0.0015$ \bic{}.

\leadsto{Two questions fall out. Those same level estimators behave very
differently on the two candidate models, which decides which model ships
(\S\ref{sec:ridge}). And 63 being wrong raises what is right, which needs a
criterion first (\S\ref{sec:criterion}).}

\subsection{Why the model is a Ridge}
\label{sec:ridge}
\prov{one\_model.py, level\_refit.py (both -{}-model \{tree,ridge\}), detrend\_sweep.py}{dev + valid}

Both models order returns about equally on dev: \bic{} 0.1060 for the Ridge against
0.1011 for the tree (\S\ref{sec:signalkind}). Running the same diagnostics on both,
from one script each with a \texttt{-{}-model} switch, holds everything but the fit
constant: the same splits, the same fold grid at embargo 0, both fitted fresh.

Hold the mapping fixed and change only how the signal is produced, over identical
valid rows. \texttt{p\_window} is the within-block rank, which is look-ahead and
immune to the block's level by construction, so it bounds what any causal rule could
reach.

\begin{center}\small
\begin{tabular}{lrrrr}
\toprule
signal produced by & tree, causal & tree, ceiling & Ridge, causal & Ridge, ceiling \\
\midrule
refit every 84 rows, 756-row window & \textbf{0.0062} & 0.0468 & \textbf{0.0984} & 0.1026 \\
1 fit, 756-row window & 0.0399 & 0.0780 & 0.0843 & 0.0861 \\
1 fit, $\sim$6600-row window & \textbf{0.0668} & 0.0896 & \textbf{0.0425} & 0.0835 \\
\bottomrule
\end{tabular}
\end{center}

The two ladders run in opposite directions. The tree's causal ordering improves
tenfold as refits are removed; the Ridge's is best with the most refits and decays
as they are taken away.

Under the shipped protocol the tree sits at 0.0062 against a ceiling of 0.0468, a
gap of 0.041, while the Ridge sits at 0.0984 against 0.1026, a gap of 0.004. The
tree's ceiling is still there, so its ordering ability survives and the loss is in
the level. The decisive number is simpler than the ratio of the gaps: the Ridge's
causal ordering, 0.0984, is more than double the tree's look-ahead ceiling of
0.0468. A perfect level repair would still leave the tree behind where the Ridge
already stands without one.

The persistence contrast agrees. Correlating a block's realised level with the
previous block's, within one fit and across a refit:

\begin{center}\small
\begin{tabular}{lrr}
\toprule
 & refit every 84 rows & 1 fit, 756-row window \\
\midrule
tree  & $\mathbf{-0.39}$ & $\mathbf{+0.65}$ \\
Ridge & $-0.02$ & $+0.12$ \\
\bottomrule
\end{tabular}
\end{center}

Same rows, same period, same algorithm; the only difference is whether a refit sat
between the two blocks. So does the table of level estimators (dev, 69 folds,
standard error 0.120):

\begin{center}\small
\begin{tabular}{lrr}
\toprule
estimator & LightGBM & Ridge \\
\midrule
the model predicting its own last 84 training rows & $+0.118$ & $\mathbf{+0.342}$ \\
the model predicting its whole training set & $+0.211$ & $\mathbf{+0.415}$ \\
mean of the training labels & $+0.212$ & $+0.366$ \\
\bottomrule
\end{tabular}
\end{center}

Nothing clears two standard errors for the tree: its level appears when the model
meets the test period and is unreadable beforehand. The same readings correlate for
the Ridge. So no estimate can claim the tree's 0.041 gap, because the level it
would have to predict is not readable in advance.

Because the level belongs to the fit, nothing chosen on one span survives to the
next. Correlating \vsbh{} across the mapping search's 40 cells, dev to valid:

\begin{center}\small
\begin{tabular}{lrr}
\toprule
 & Pearson & Spearman \\
\midrule
tree & $\mathbf{-0.519}$ & $-0.506$ \\
Ridge & $\mathbf{+0.426}$ & $+0.376$ \\
\bottomrule
\end{tabular}
\end{center}

Choosing a configuration on dev actively hurts the tree on valid. This project's
own historical benchmark for the same number is $+0.070$, across 37
volatility-overlay candidates.

The consequences show up under the identical shipped mapping:

\begin{center}\small
\begin{tabular}{lrr}
\toprule
 & dev & valid \\
\midrule
tree & $+0.089$ & $\mathbf{-0.069}$ \\
Ridge & $+0.034$ & $\mathbf{+0.043}$ \\
\bottomrule
\end{tabular}
\end{center}

and in the search of \S\ref{sec:control}, where configurations clearing dev and
valid come to 0 of 210 for the tree against 90 of 210 for the Ridge. The tree
clears nothing anywhere in that grid, under any turnover control.

Ship the Ridge. One mechanism explains all three measurements: each of the tree's
fits carries its own output level, a level that belongs to a fit cannot survive
being refitted, and so nothing selected on one span carries to the next. The
Ridge's output level barely moves between fits, which is why its configuration
space transfers and the tree's inverts.

That also settles \S\ref{sec:signalkind}'s debt. The tree was owed a magnitude
encoding, since a rank-only mapping is lossless for the Ridge ($t$ on the magnitude
$= 0.74$) and lossy for the tree ($t = 3.52$). The tree is dropped because nothing
selected on it transfers, which is a separate finding from the encoding it never
got.

\paragraph{A limit on this section.} \texttt{valid} is ten blocks, and the tree arm
is seven bags at one seed against the Ridge's closed form. The directions agree
across three independent measurements; the individual values are not precise.

\leadsto{The model is settled. \S\ref{sec:reference} found that removing nothing
orders best, so the criterion has to be checked before the length of what to remove
can be swept.}

\subsection{Select on the scored metric, not on the ordering}
\label{sec:criterion}
\prov{detrend\_sweep.py -{}-grid -{}-axis}{Ridge, dev + valid}

\begin{center}\small
\begin{tabular}{llrrrrr}
\toprule
span & detrend & \bic{} & \vsbh{} & exposure & vol ratio & turnover \\
\midrule
dev & none & \textbf{0.0950} & $+0.028$ & 1.004 & 1.211 & 0.191 \\
dev & roll21 & 0.0558 & $\mathbf{+0.060}$ & 1.005 & 1.195 & 0.298 \\
valid & none & \textbf{0.089} & $\mathbf{-0.067}$ & 1.006 & 1.177 & 0.124 \\
valid & roll8 & 0.082 & $\mathbf{+0.384}$ & 0.999 & 1.143 & 0.344 \\
valid & roll63 & 0.064 & $-0.072$ & 0.989 & 1.160 & 0.131 \\
\bottomrule
\end{tabular}
\end{center}

The configuration that orders best pays least, on both spans. On valid the two
rules have nearly the same ordering, 0.089 against 0.082, the same average
exposure, 1.006 against 0.999, and the same realised volatility, 1.177 against
1.143. They differ by 0.45 in adjusted Sharpe. Ordering skill and risk taken are
both ruled out by those columns, which leaves when the position moves. \bic{} is a
rank correlation weighting every day equally, while the metric is carried by a
handful of large days, and a rank correlation cannot see the difference.

So \bic{} drops to a diagnostic and configurations get chosen on adjusted Sharpe
against buy-and-hold. The responsiveness account is inferred from what the other
columns exclude rather than measured; a return-weighted IC, or the same comparison
split by the size of the day's move, would test it directly.

\leadsto{With the criterion fixed, the length of the removed mean can be swept on
the metric.}

\subsection{The length of the mean that gets removed}
\label{sec:meanlen}
\prov{mapping\_form.py -{}-axes, detrend\_sweep.py -{}-axis}{dev and valid}

Dev alone does not resolve this. \texttt{mapping\_form.py -{}-axes}, pooled on dev,
puts \texttt{mu\_window} $=21$ at the top of every rank-window column, but the
paired contrast against 63 is $+0.0053$ against a \detec{} of 0.106, twenty times
inside the noise floor. Valid can separate the short lengths.

\begin{center}\footnotesize
\begin{tabular}{lrrrrrrrrrr}
\toprule
length & 5 & 8 & 13 & 21 & 34 & 42 & 63 & 126 & 252 & none \\
\midrule
\vsbh{} & $+0.264$ & $\mathbf{+0.384}$ & $+0.198$ & $+0.031$ & $\mathbf{-0.025}$ & $-0.010$ & $-0.072$ & $-0.082$ & $-0.099$ & $-0.067$ \\
turnover & 0.399 & 0.344 & 0.275 & 0.217 & 0.171 & --- & 0.131 & --- & 0.129 & 0.124 \\
\bottomrule
\end{tabular}
\end{center}

The break sits between 21 and 34. Block-aligned alternatives in the same grid fail
on valid in every variant:

\begin{center}\small
\begin{tabular}{lrr}
\toprule
detrend & dev & valid \\
\midrule
none at all & $+0.028$ & $-0.067$ \\
mean of the previous 1 complete block & $-0.010$ & $-0.074$ \\
previous 3 blocks & $+0.003$ & $-0.096$ \\
previous 6 blocks & $-0.008$ & $-0.062$ \\
every previous block & $+0.017$ & $-0.067$ \\
\textbf{rolling 8} & $\mathbf{+0.074}$ & $\mathbf{+0.384}$ \\
\bottomrule
\end{tabular}
\end{center}

A trailing mean of 5 to 21 rows. Removing nothing loses money despite ordering best
(\S\ref{sec:criterion}); removing a long mean loses money and orders badly; the
short end is the only range that does both acceptably. Lengths $\ge 34$ are out, so
is removing nothing, and so are the block-aligned means. A level decided at refit
time and held constant through the block never straddles a boundary, and it is
negative on valid in every variant, so the answer is neither ``avoid the block
boundary'' nor ``leave the level alone''.

The mechanism stays open. A clean account exists for a model whose level belongs to
the fit: a mean of length $L$ reaches into the previous fit on about $L/84$ of days,
and consecutive fits' levels are uncorrelated. \S\ref{sec:ridge} shows the Ridge is
not such a model, so the account does not apply. What is established is the
empirical break, not why it sits at 21.

\leadsto{The turnover row above is the price of the short end. The next step decides
whether to pay it.}

\subsection{Turnover is the cost, and the metric does not charge it}
\label{sec:turnover}
\prov{detrend\_sweep.py -{}-axis}{Ridge, valid}

Turnover is monotone in the length of the mean, from 0.399 at length 5 to 0.124 at
no removal. A shorter mean means a faster-moving position, through the same
operation. At the short end the rule changes 34 to 40\% of notional every day, and
unconstrained configurations across the whole grid run from 25 to 50\%.

The metric contains no execution. It scores a position series, so a rule trading
40\% of notional a day and one trading 5\% receive identical treatment. Nothing in
the metric constrains turnover, and nothing in dev or valid does either.

Trading 40\% of notional a day is not a deployable product whatever the metric
says, so the position is required to be smoothed and the parameters are chosen
inside that constraint. This is a prior about what kind of rule is worth having.
The data does not supply it, and it rules out the unsmoothed family on grounds the
data does not support.

\leadsto{Given that turnover has to come down, the question is which way of
bringing it down destroys the least edge.}

\subsection{Continuous smoothing gives up the least edge per unit of turnover}
\label{sec:control}
\prov{detrend\_sweep.py -{}-control}{Ridge, dev + valid}

15 ways of computing the removed mean $\times$ 7 controls $\times$ 2 rank windows,
the controls applied to the position after the scale is solved, read as edge
retained per unit of turnover spent and averaged over all fifteen means.

\begin{center}\small
\begin{tabular}{lrrr}
\toprule
control (dev) & \vsbh{} & turnover & edge per unit turnover \\
\midrule
unconstrained & 0.059 & 0.324 & 0.18 \\
band 0.05 & 0.044 & 0.269 & 0.16 \\
band 0.10 & 0.030 & 0.224 & 0.13 \\
band 0.20 & 0.000 & 0.154 & \textbf{0.00} \\
EWMA half-life 2 & 0.027 & 0.095 & 0.28 \\
EWMA half-life 3 & 0.031 & 0.070 & 0.44 \\
\textbf{EWMA half-life 5} & 0.035 & 0.047 & \textbf{0.74} \\
\bottomrule
\end{tabular}
\end{center}

\begin{center}\small
\begin{tabular}{lrrr}
\toprule
clearing dev and valid, by family & unconstrained & band & EWMA \\
\midrule
short mean (3--21) & 20 of 22 & 41 of 66 & \textbf{29 of 66} \\
long mean (34--63) & \textbf{0 of 6} & \textbf{0 of 18} & \textbf{0 of 18} \\
no mean removed & \textbf{0 of 2} & \textbf{0 of 6} & \textbf{0 of 6} \\
\bottomrule
\end{tabular}
\end{center}

A no-trade band is dominated at every level. Band 0.20 gives up the entire edge,
$+0.000$, to halve turnover, while an EWMA of half-life 5 keeps 59\% of it at a
seventh of the trading. What each does to the level explains why: a band holds the
position unchanged on most days and loses track of the level, while an EWMA moves
every day by less and keeps it. The control is an EWMA on the position.

The family separation in the second table is not a turnover effect. The long-mean
and no-removal families clear nothing under any control, including no control at
all, so \S\ref{sec:meanlen}'s result holds however turnover is handled.

\leadsto{Every parameter now has a value or a range. What remains is choosing among
the survivors by a rule fixed in advance.}

\section{Selection and result}

\subsection{Selection}
\label{sec:selection}
\prov{detrend\_sweep.py -{}-control}{dev + valid only, scored by the competition metric with nothing added}

The smoothing constraint from \S\ref{sec:turnover} leaves 90 configurations: 15 ways
of computing the removed mean $\times$ 3 EWMA half-lives $\times$ 2 rank windows.
The pre-registered rule is maximin, maximise the smaller of dev and valid, which
requires both spans to stand up and needs no threshold.

\begin{center}\small
\begin{tabular}{llrrrr}
\toprule
removed mean & EWMA half-life & rank window & dev & valid & maximin \\
\midrule
\textbf{ewm3} & \textbf{3} & \textbf{252} & $+0.036$ & $+0.063$ & \textbf{0.036} \\
ewm3 & 3 & 504 & $+0.035$ & $+0.064$ & 0.035 \\
mean8 & 5 & 252 & $+0.034$ & $+0.043$ & 0.034 \\
med8 & 5 & 252 & $+0.034$ & $+0.068$ & 0.034 \\
ewm3 & 2 & 252 & $+0.033$ & $+0.118$ & 0.033 \\
mean5 & 2 & 504 & $+0.033$ & $+0.216$ & 0.033 \\
ewm3 & 2 & 504 & $+0.032$ & $+0.122$ & 0.032 \\
mean5 & 3 & 504 & $+0.032$ & $+0.154$ & 0.032 \\
\bottomrule
\end{tabular}
\end{center}

The top of that ranking is a tie. The first eight rows differ by 0.004 in maximin
score, far inside any resolution these spans support, and 29 of the 90 clear both.
So three rules are applied and all three are carried forward.

\begin{center}\small
\begin{tabularx}{\textwidth}{ll>{\raggedright\arraybackslash}X}
\toprule
rule & selects & why this rule \\
\midrule
A, maximise $\min(\text{dev},\text{valid})$ &
\texttt{ewm3} $\times$ hl 3 $\times$ rank 252 &
pre-registered; requires both spans to hold, needs no threshold \\[4pt]
B, maximise valid subject to dev $>0$ &
\texttt{mean5} $\times$ hl 2 $\times$ rank 504 &
uses valid for what valid is for, with dev only as a gate \\[4pt]
C, maximise $\mathrm{mean}(\text{dev},\text{valid})$ &
\texttt{mean5} $\times$ hl 2 $\times$ rank 504 &
equal weight to both spans; picks the same cell as B \\[4pt]
D, best flat-window mean &
\texttt{mean8} $\times$ hl 5 $\times$ rank 252 &
the three ways of computing the mean tie at the top, and a flat window is the
plainest form of the operation \\
\bottomrule
\end{tabularx}
\end{center}

B and C agree, so three rules give two configurations. D is added because the choice
between an exponentially weighted, a flat-window and a median mean is a tie on the
evidence, and a reader reproducing this should see the plainest version. A reported
result that depends on which rule was used is worth less than one that does not, so
all of them are reported.

\leadsto{The rule is now fully specified. What follows is its statement, what the
data actually fixed, and the one reading of \texttt{test}.}

% Keep the heading and the whole code block on one page: if less than 21 lines
% are left, start them on the next page. Raise the number if the block grows.
\Needspace{21\baselineskip}
\subsection{The rule}
\label{sec:rule}
\prov{pipeline.py}{all three spans}

\begin{lstlisting}
Ridge(alpha=1000)  |  rolling 756  |  refit every 84  |  1132 columns
    preprocessing: ffill -> standardise (degenerate columns pinned)
                   -> clip to +/-10 training SD
    |  s(t)
v(t) = s(t) - recent mean of s             remove the slow level, once
    |
p(t) = percentile of v(t) in a trailing window    turn it into an order
    |
z(t) = sqrt(3) * (2*p(t) - 1)
    |
tau(t) solved daily on trailing 756 so realised vol ratio = 1.2
    |
w*(t) = clip(1 + tau(t)*z(t), 0, 2)
    |
w(t)  = EWMA(w*, a few days)               hold turnover near 10% a day
\end{lstlisting}

Three settings of the two free lengths and the rank window, one per rule:

\begin{center}\small
\begin{tabular}{llrl}
\toprule
rule & removed mean & rank window & position EWMA \\
\midrule
A, maximin & exponential, half-life 3 & 252 & half-life 3 \\
B/C, max valid and max mean & flat window of 5 & 504 & half-life 2 \\
D, flat-window mean & flat window of 8 & 252 & half-life 5 \\
\bottomrule
\end{tabular}
\end{center}

\paragraph{Which parameters the data actually fixed.}

\begin{center}\small
\begin{tabular}{p{0.21\textwidth}p{0.13\textwidth}p{0.40\textwidth}l}
\toprule
parameter & value & separable & where \\
\midrule
model & Ridge & yes, transfer dev$\to$valid $+0.43$ against $-0.52$ & \S\ref{sec:ridge} \\
curve shape & linear & \textbf{no}, 7 forms inside 0.021 & \S\ref{sec:curve} \\
mean exposure $\bar w$ & 1.00 & \textbf{no}, flat over 0.9--1.1 & \S\ref{sec:amplitude} \\
volatility budget & 1.2 & yes & \S\ref{sec:amplitude} \\
scale re-solve cadence & daily & pooled yes ($+0.049$), per-fold no & \S\ref{sec:amplitude} \\
the removed mean & short, 3--8 & family yes ($\le 21$ against $\ge 34$), within-family \textbf{no} & \S\ref{sec:reference}--\S\ref{sec:meanlen} \\
rank window & 252 or 504 & \textbf{no}, both clear & \S\ref{sec:meanlen}, \S\ref{sec:selection} \\
turnover control & EWMA, hl 2--5 & control \textbf{yes}, half-life \textbf{no} & \S\ref{sec:turnover}--\S\ref{sec:control} \\
\bottomrule
\end{tabular}
\end{center}

Half of these are ties, which is why \S\ref{sec:selection} carries three
configurations. The curve's shape, the mean exposure, the rank window, the length
inside the short family and the half-life that smooths the position are all ``any
value in this range''. What the data fixed is the model, the volatility budget, the
solve cadence, the family the removed mean belongs to, and that the turnover control
is an EWMA rather than a band.

\subsection{Result}
\label{sec:result}

Every row was selected on dev and valid alone, by the correspondingly lettered rule.
\texttt{test} was scored afterwards. Buy-and-hold is $0.415 / 0.968 / 0.302$.

\begin{center}\small
\begin{tabular}{llrrr}
\toprule
rule & configuration & dev & valid & test \\
\midrule
A, maximin & \texttt{ewm3} $\times$ hl3 $\times$ rank252 & 0.451 & 1.030 & \textbf{0.435} \\
B/C, max valid and max mean & \texttt{mean5} $\times$ hl2 $\times$ rank504 & 0.448 & 1.183 & \textbf{0.448} \\
D, flat-window mean & \texttt{mean8} $\times$ hl5 $\times$ rank252 & 0.449 & 1.011 & \textbf{0.399} \\
\bottomrule
\end{tabular}
\end{center}

\begin{center}\small
\begin{tabular}{lrrrrr}
\toprule
against buy-and-hold & dev & valid & test & turnover (test) & vol ratio (test) \\
\midrule
A & $+0.036$ & $+0.063$ & $\mathbf{+0.133}$ & 0.106 & 1.057 \\
B/C & $+0.033$ & $+0.216$ & $\mathbf{+0.146}$ & 0.151 & 1.077 \\
D & $+0.034$ & $+0.043$ & $\mathbf{+0.096}$ & 0.068 & 1.041 \\
\bottomrule
\end{tabular}
\end{center}

All three clear buy-and-hold on all three spans. Average exposure is 1.00
everywhere, turnover runs 4--15\% of notional a day, and both penalty factors come
out at exactly 1.000, so realised volatility never approaches the 1.2 kink and the
strategy never underperforms the market.

The strongest statement the evidence supports is about the family rather than any
one row. Of the 90 smoothed configurations, 29 clear both dev and valid, and all 29
are positive on test: mean $+0.129$, minimum $+0.079$, maximum $+0.168$. The
long-mean and no-removal families clear nothing under any turnover control,
including none at all.

Three checks made on dev and valid also held when test was read, and none of them
entered a decision. The length sweep of \S\ref{sec:meanlen} keeps its sign in 11 of
12 cells. The cell-to-cell transfer of \S\ref{sec:ridge} stays positive for the
Ridge, $+0.224$ from dev and $+0.476$ from valid, and stays negative for the tree,
$-0.741$ from valid. Under the shipped mapping the tree reaches $+0.012$ against the
Ridge's $+0.096$, so the model choice would have been the same had test been
available.

The smoothing constraint cost measured score. Under this metric more turnover scores
better, and averaged over the grid an unconstrained position beats every smoothed
one on both development spans (\S\ref{sec:turnover}). Requiring the position to be
smoothed therefore gave up score on dev and valid, on deployability grounds.

\subsection{What this is not}
\label{sec:limits}

The edge is small: $+0.036$ on dev against a pre-registered threshold of $+0.10$. It
clears buy-and-hold on all three spans for a reason that replicates across a family
of configurations, and it is scored on daily equity-index exposure by a metric that
charges nothing for trading.

Two mechanisms stay open: why a short mean pays when it orders no better than
removing nothing (\S\ref{sec:criterion}), and why the break sits at 21
(\S\ref{sec:meanlen}). The account that fitted the LightGBM does not apply to the
model that ships.

\appendix
\section{The volatility budget is left unspent on purpose}
\label{sec:budget}
\prov{scale\_order.py}{dev + valid}

Smoothing comes after the scale is solved, so realised volatility falls from the
1.19 the solver aims at to about 1.06. The penalty never binds and part of the
allowance goes unused, which raises whether to reclaim it.

An EWMA is linear, so
$\mathrm{EWMA}(\bar w + \tau z) = \bar w + \tau\,\mathrm{EWMA}(z)$. Solving the
scale on the smoothed signal aims it at what will actually be held. Three arms, the
third a control that filters the signal but keeps the old scale.

\begin{center}\small
\begin{tabular}{llrrrrr}
\toprule
config & arm & dev & valid & maximin & turnover (dev) & vol ratio \\
\midrule
A, maximin & smooth after the solve (shipped) & $\mathbf{+0.036}$ & $+0.063$ & \textbf{0.036} & 0.087 & 1.083 \\
 & solve on the smoothed signal & $+0.019$ & $+0.093$ & 0.019 & 0.157 & 1.190 \\
 & filter the signal, keep the old scale & $+0.036$ & $+0.067$ & 0.036 & 0.087 & 1.080 \\
B/C, max valid & smooth after the solve (shipped) & $\mathbf{+0.033}$ & $+0.216$ & \textbf{0.033} & 0.131 & 1.090 \\
 & solve on the smoothed signal & $+0.008$ & $+0.343$ & 0.008 & 0.222 & 1.188 \\
 & filter the signal, keep the old scale & $+0.033$ & $+0.220$ & 0.033 & 0.131 & 1.088 \\
D, flat-window & smooth after the solve (shipped) & $\mathbf{+0.034}$ & $+0.043$ & \textbf{0.034} & 0.057 & 1.061 \\
 & solve on the smoothed signal & $+0.019$ & $+0.068$ & 0.019 & 0.129 & 1.191 \\
 & filter the signal, keep the old scale & $+0.034$ & $+0.047$ & 0.034 & 0.057 & 1.057 \\
\bottomrule
\end{tabular}
\end{center}

The algebra holds. The third arm matches the first to three decimals on dev in every
configuration, and the clip binds on 4--5\% of days under the re-solved arm and on
none under the shipped one. Re-solving does push the realised ratio from about
1.06--1.09 up to 1.19, and it cuts the maximin score by half or more in all three
configurations while roughly doubling turnover. It loses under the same rule
everything else was chosen by.

This is \S\ref{sec:amplitude} from the other side. The unpenalised Sharpe saturates
near a ratio of 1.3, so a wider swing amplifies whatever is in the signal, including
the noise. The gain sits entirely on valid, where the signal is strong, and the loss
sits on dev's 69 blocks. The unspent allowance is not free money and the shipped
order stands.

% Keep this short appendix together, so its last lines never end up alone on a
% page of their own.
\Needspace{14\baselineskip}
\section{Reproducing}
\label{sec:repro}

Every number here is written to \texttt{kaggle\_hull/probe/} by the script named in
its step, and re-reading it needs no refitting: the signal caches are keyed by
feature-table row, so a later stage replays without re-running an earlier one.

\begin{lstlisting}[language=bash]
python stage1/Hull_Tactical_feature_engineering.py   # the 1132 columns
python tools/split_data.py --commit                  # the three slices
python tools/audit_config.py                         # exits with the mismatch count
python stage2/pipeline.py --commit                   # the result
\end{lstlisting}

\texttt{docs/RUNBOOK.md} gives all eighteen commands in order with their costs.
\texttt{detrend\_sweep.py} and \texttt{scale\_order.py} score dev and valid only;
\texttt{-{}-reveal-test} opens the third span.

\end{document}
